\section{The Hamiltonian and its auxiliary twists}\label{sec:model}

The argument uses the same partition function in two different ways: as
a Gibbs trace in the physical Hilbert space and as a moment of a spatial
operator. We first choose coordinates that make both descriptions explicit.
We also establish an elementary energy bound used to control the finite
exponential approximations.

\subsection{Spin coordinates}

We use the self-adjoint spin-one matrices. This convention does not restrict
the spectrum in the formulation where self-adjointness is not assumed.

\begin{lemma}[Choice of representation]\label{lem:representation}
Let $S^x,S^y,S^z\in M_3(\C)$ satisfy
\[
 [S^\alpha,S^\beta]
   =i\sum_\gamma\varepsilon_{\alpha\beta\gamma}S^\gamma,
 \qquad \sum_\alpha(S^\alpha)^2=2I.
\]
They are simultaneously similar to the standard self-adjoint spin-one
matrices. If the original matrices are self-adjoint, the similarity can be
chosen unitary. The corresponding tensor-power similarity preserves the
complete spectrum of every periodic Hamiltonian $H_n$.
\end{lemma}

\begin{proof}
Put $S^\pm=S^x\pm iS^y$. The assumptions give
\begin{equation}\label{eq:ladder-identities}
 [S^z,S^\pm]=\pm S^\pm,\qquad
 S^-S^+=2I-(S^z)^2-S^z,\qquad
 S^+S^-=2I-(S^z)^2+S^z.
\end{equation}
Choose an $S^z$ eigenvector $v$ whose eigenvalue $m$ has maximal real part.
Then $S^+v=0$, since otherwise it has eigenvalue $m+1$. Thus $m=1$ or
$m=-2$. The latter case would permit indefinite nonzero lowering: a last
nonzero lowered vector would have weight $-2-k$, whereas
$S^+S^-w=0$ forces its weight to be $-1$ or $2$. This contradicts finite
dimension. Hence $m=1$.

The vectors
\[
 v,\qquad \frac{S^-v}{\sqrt2},\qquad \frac{(S^-)^2v}{2}
\]
are nonzero by \eqref{eq:ladder-identities} and have distinct weights
$1,0,-1$. They form a basis. There is no further lowered vector, because
the basis has no weight $-2$. The identities give raising and lowering
coefficients $\sqrt2$, so this is the standard spin-one representation.
If the original matrices are self-adjoint and $\norm v=1$, the three
vectors are orthonormal. Finally, conjugating each tensor factor conjugates
every term of $H_n$.
\end{proof}

In the weight basis $|s\rangle$, $s\in\{-1,0,1\}$, the diagonal generator
is $S^z|s\rangle=s|s\rangle$ and the nonzero ladder entries are $\sqrt2$.
We will also use the Cartesian basis, indexed by the three axes, in which
\[
 (S^\alpha)_{ij}=-i\varepsilon_{\alpha i j},\qquad
 G_\alpha=iS^\alpha.
\]
These matrices satisfy the same relations. Each $G_\alpha$ is real and
has operator norm one. For the two-site interaction
$h=\sum_\alpha S^\alpha\otimes S^\alpha$, contraction of the epsilon
tensors gives
\begin{equation}\label{eq:cartesian-bond}
 -h=\sum_\alpha G_\alpha\otimes G_\alpha,
 \qquad h=F-|\Omega\rangle\langle\Omega|,
 \qquad \Omega=\sum_{i=1}^3|ii\rangle,
\end{equation}
where $F(u\otimes v)=v\otimes u$. In particular, $h$ has eigenvalues
$-2,-1,1$, and $h\le I$.

\subsection{Twists and partition functions}

Number the sites $0,\ldots,n-1$, with $n\ge4$. A proper Cartesian rotation
$g\in SO(3)$ is unitary on $\C^3$. Define
\begin{equation}\label{eq:twisted-hamiltonian}
 H_n(g)=\sum_{j=0}^{n-2}h_{j,j+1}+h_{n-1,0}(g),
 \qquad
 h_{n-1,0}(g)=\sum_\alpha
 S_{n-1}^\alpha(gS^\alpha g^{-1})_0.
\end{equation}
Thus only the last bond is conjugated, on its site-$0$ factor. The
Hamiltonian is self-adjoint, and $H_n(I)=H_n$. Fix
\begin{equation}\label{eq:shift}
 a=\frac{700741}{500000},\qquad
 Z_n(b,g)=\Tr\exp[-b(H_n(g)+anI)],\qquad Z_n(b)=Z_n(b,I)
 \quad (b>0).
\end{equation}
All these partition functions are strictly positive. The scalar shift
will cancel from the purity ratios and from the final spectral gap.
With this normalization, the trial bounds in Lemma~\ref{lem:trial-inputs}
give $E_0(n)+an<0$ at $n=72,120$, and hence $Z_n(b)>1$ for every $b>0$.
These lower bounds are paired with finite thermal upper bounds in
Section~\ref{sec:initialization}.

We need the diagonal rotations
\[
 P_x=\operatorname{diag}(1,-1,-1),\quad
 P_y=\operatorname{diag}(-1,1,-1),\quad
 P_z=\operatorname{diag}(-1,-1,1),\qquad
 D_2=\{I,P_x,P_y,P_z\},
\]
and the cyclic rotation $C:e_x\mapsto e_y\mapsto e_z\mapsto e_x$.
Write $P=P_x$. The three nonidentity elements of $D_2$ have equal
partition functions, by a global cyclic rotation of all sites. The
rotation $C$ has angle $2\pi/3$ about the diagonal axis. After a global
rotation and a change to the weight basis, $P$ and $C$ are respectively
represented by $\exp(-i\theta S^z)$ with $\theta=\pi$ and
$\theta=2\pi/3$.

For later comparison with finite matrices, a word $x=(x_0,\ldots,x_{n-1})$
in the weight basis has diagonal entry $\sum_j x_{j-1}x_j$, with indices
modulo $n$. Each permitted step changing two neighboring letters by
opposite units has coefficient one. For a seam step from column $y$ to
row $x$, this coefficient is
\begin{equation}\label{eq:seam-phase}
 \exp\bigl(i\theta(y_0-x_0)\bigr).
\end{equation}
Indeed $gS^\pm g^{-1}=e^{\mp i\theta}S^\pm$. These formulas identify the
finite matrices used below with the physical Hamiltonians in
\eqref{eq:twisted-hamiltonian}.

\subsection{A uniform energy bound}

\begin{lemma}[Energy bounds]\label{lem:energy}
For every $n\ge4$ and every $g\in SO(3)$,
\begin{equation}\label{eq:energy-bounds}
 -\frac{3n}{2}I\le H_n(g)\le nI.
\end{equation}
\end{lemma}

\begin{proof}
The upper bound follows from $h\le I$ and unitary conjugation of the
seam. For the lower bound, consider $T=-h_{12}-h_{23}$ on three open
sites in Cartesian coordinates. By \eqref{eq:cartesian-bond}, an unequal
pair swaps with coefficient $-1$, and an equal pair changes to either of
the other equal pairs with coefficient $1$. There are no diagonal terms.

Assign a positive weight $q$ to each three-letter word. The following
table lists every equality pattern; $a,b,c$ denote distinct letters.
The last column sums the weights of the neighbors of a row word, with
the absolute matrix coefficients.
\begin{center}
\begin{tabular}{c|c|c}
word type & $q$ & absolute weighted row sum\\ \hline
$aaa$ & $4$ & $3+3+3+3=12$\\
$aab$ or $abb$ & $3$ & $4+3+2=9$\\
$aba$ & $2$ & $3+3=6$\\
$abc$ & $2$ & $2+2=4$
\end{tabular}
\end{center}
Every row sum is at most $3q$. If $D$ is the diagonal matrix of these
weights, $D^{-1}TD$ has absolute row sums at most three. Its spectral
radius is therefore at most three. Since $T$ is self-adjoint and similar
to this matrix, $h_{12}+h_{23}\ge-3I$.

An open triad containing the seam is unitarily equivalent to an
untwisted triad: delete the seam bond and rotate all sites in one of
the two resulting components. The internal interaction of that component
is rotation invariant. Thus the same lower bound holds for every
consecutive pair of bonds in $H_n(g)$. Summing the $n$ triad inequalities
counts each bond twice and proves \eqref{eq:energy-bounds}.
\end{proof}
